% LaTeX resume using res.cls
\documentclass[margin]{res}
\usepackage[T1]{fontenc}
\usepackage[utf8]{inputenc}
%\usepackage{helvetica} % uses helvetica postscript font (download helvetica.sty)
%\usepackage{newcent}   % uses new century schoolbook postscript font 
\setlength{\textwidth}{5.1in} % set width of text portion
\setlength{\emergencystretch}{1.5em}
\usepackage[protrusion=true,expansion=false]{microtype}
\usepackage{needspace}
\usepackage{hyperref}
\usepackage{xcolor}
\hypersetup{colorlinks,urlcolor=blue}
\begin{document}
	
	% Center the name over the entire width of resume:
	\moveleft.5\hoffset\centerline{\Large\bf Jan Schulz-Gebhard}
	% Draw a horizontal line the whole width of resume:
	\moveleft\hoffset\vbox{\hrule width\resumewidth height 1pt}\smallskip
	% address begins here
	% Again, the address lines must be centered over entire width of resume:
	\moveleft.5\hoffset\leftline{University of Bamberg}% \hfill Phone: (801) 581-7481}
\moveleft.5\hoffset\leftline{Institute of Economics\hfill \href{mailto:jan.schulz-gebhard@uni-bamberg.de}{\texttt{jan.schulz-gebhard@uni-bamberg.de}}}
\moveleft.5\hoffset\leftline{\hfill \href{https://schulz-gebhard.github.io/}{\texttt{schulz-gebhard.github.io}}}
\moveleft.5\hoffset\leftline{Feldkirchenstr. 21}
\moveleft.5\hoffset\leftline{96052 Bamberg}
\moveleft.5\hoffset\leftline{Room F21/02.35}
\moveleft.5\hoffset\leftline{Chair of International Economics}

\vspace{.3cm}
\moveleft.5\hoffset\leftline{* April 5, 1991, Heilbronn (Germany), three children}
\begin{resume}
	\setlength{\parskip}{0.52\baselineskip}
	
	
	
	\section{Education}
	{\sl PhD} University of Bamberg \hfill March 2023 \\
	Dissertation: \textit{``Essays on Aggregation With Heterogeneous and Interacting Agents''} (summa cum laude)\\
	Committee: M. Milakovi\'c (supervisor), F. Westerhoff, C. Proa\~no
	
	
	{\sl PhD Scholarship} Hans-B\"ockler-Stiftung \hfill 2018 -- 2020\\
	Bamberg Doctoral Research Group on Behavioral Macroeconomics 
	
	{\sl Master of Science} European Economic Studies \hfill 2017\\
	University of Bamberg\\ 
	%	Final Grade: 1.4\\
	Thesis: \textit{An Empirical Investigation Into the German Wealth Distribution}\\
	%	(Thesis Grade: 1.0)\\
	Supervisor: Professor Mishael Milakovi\'c, PhD 
	
	{\sl Bachelor of Arts} Philosophy \& Economics \hfill 2015\\ 
	University of Bayreuth\\
	Thesis: \textit{The Impact of Nonhomothetic Preferences and the Threat of Parallel Trade: An Analysis of Trade Patterns under Monopolistic Competition}\\ %(Thesis Grade 1.3)\\
	%Final Grade: 1.6\\
	Supervisor:  Professor Hartmut Egger
	
	{\sl Semester Abroad} Economics and Economic History \hfill 2013\\
	University of Graz\\ 
	
	\section{Academic Appointments}
	\textbf{Postdoctoral Researcher}
	Within ANR-DFG funded project ``The Metabolism of Capital: A Historical Account and Statistical Analysis of Corporate Profitability in France and Germany'' (MECAPHIST) (2024 -- 2027).
	
	\textbf{Postdoctoral Researcher and Lecturer}, Chair of International Economics, University of Bamberg, Professor Mishael Milakovi\'c, PhD (since 2023).
	
	\textbf{Research Assistant and Junior Lecturer}, Chair of International Economics, University of Bamberg, Professor Mishael Milakovi\'c, PhD (2018 -- 2023).
	
	\textbf{Student Assistant}, Chair of International Economics, University of Bamberg, Professor Mishael Milakovi\'c, PhD (2015 -- 2017).

	\textbf{Teaching Assistant}, Lecture Social Choice Theory, University of Bayreuth, Dr. Gordian Haas (2014).
	
	\textbf{Research Assistant}, Institute for Middle Eastern Democracy, Berlin, Saba Farzan (2014).
	
	\textbf{Teaching Assistant}, Lecture Introduction to Philosophy, University of Bayreuth, Professor Allard Tamminga (2011 -- 2012).\\
	
	\enlargethispage{3\baselineskip}
	\section{Parental Leave}
	\textbf{12 months total, 2023 -- 2026} April, August--September 2023; February--April 2024; June, October--November 2025; February--April 2026.\\
	
	\section{Research Identity}
	Mathematical, quantitative and computational social scientist studying the determinants and consequences of networked heterogeneity, with particular emphasis on how economic inequality and socially structured information environments shape public perceptions, redistribution preferences and political behaviour. Trained in economics, I combine social-network analysis, surveys and experiments, formal and computational modelling, and historical and spatial microdata, with applications to economic inequality, production networks and other complex socioeconomic systems.\\
	
	\textbf{Fields:} Macroeconomics; Public Economics; Political Economy; Inequality and Redistribution; Public Opinion and Political Perception; Social Networks; Computational Social Science; Production Networks; Historical Political Economy; Intersectionality and Discrimination.\\
	
	\section{Peer-Reviewed Articles}
	\textbf{Schulz, Jan}, Agoha, Caleb, Gebhard, Anna, Gregg, Bettina and Mayerhoffer, Daniel (2026): ``Excessive white male privilege biases the measurement of intersectional wage discrimination''. \textit{Humanities and Social Sciences Communications}, 13: 631.
	
	Ipsen, Leonhard and \textbf{Schulz, Jan} (2026): ``The (Dis)Equalizing Effects of Production Networks'', \emph{Applied Economics Letters}, 33(6): 830--834.
	
	\textbf{Schulz, Jan} and Mayerhoffer, Daniel (2026): ``Consumption emulation and demand regimes: An inclusive modeling approach'', \textit{Review of Political Economy}, 38(2): 589--612.
	
	Ipsen, Leonhard, Aminian, Armin and \textbf{Schulz, Jan} (2025): ``Stress Testing Inflation Exposure: Systemically Significant Prices and Asymmetric Shock Propagation in the EU''. \textit{Structural Change and Economic Dynamics}, 74: 713--724.
	
	Mayerhoffer, Daniel and \textbf{Schulz, Jan} (2026): ``Social Segregation, Misperceptions, and Emergent Cyclical Choice Patterns'', \emph{Public Choice}, 207: 611--644.
	
	Weber, Jan D. and \textbf{Schulz, Jan} (2025): ``Growing Differently: European Integration and Regional Cohesion'', \emph{Journal of Economic Issues}, 59: 894 -- 912.
	
	\textbf{Schulz, Jan} and Mayerhoffer, Daniel (2024): ``Social Sampling, Self-Anchoring and Redistribution'', \emph{Economics Letters}, 244: 111942.
	
	
	\textbf{Schulz, Jan} and Milakovi\'c, Mishael (2023): ``How Wealthy Are the Rich?'', \textit{Review of Income and Wealth}, 69(1): 100--123.\\
	Media coverage: B\"ockler Impuls 19/2021 ``Untersch\"atzter Reichtum'', Paper of the Week at Momentum Institut, interview with stepstone.de ``Ab wann ist man reich?'' (January 09, 2024).
	
	\textbf{Schulz, Jan} and Mayerhoffer, Daniel (2023): ``A Network Approach to Expenditure Cascades'', \emph{Review of Behavioral Economics}, 10(3): 229--262.
	
	\textbf{Schulz, Jan}, Mayerhoffer, Daniel and Gebhard, Anna (2023): ``Soziale Vergleiche und wahrgenommene Ungleichheit'', \textit{Der \"Offentliche Sektor -- The Public Sector}, 48(2): 15--18. (German-language summary and outlook paper for the Egon-Matzner-Preis f\"ur Sozio\"okonomie 2022).
	
	
	\textbf{Schulz, Jan}, Mayerhoffer, Daniel and Gebhard, Anna (2022): ``A Network-Based Explanation of Inequality Perceptions'', \textit{Social Networks}, 70: 306--324.\\
	Media coverage: Podcast ``In der Wirtschaft'' (No. 61), Frankfurter Neue Presse (March 23, 2024) and M\"unchener Merkur (March 31 and July 18, 2024), Handelsblatt (April 11, July 06, and September 09, 2025).
	
	Mayerhoffer, Daniel and \textbf{Schulz, Jan} (2022): ``Perception and Privilege'', \textit{Applied Network Science}, 7(32).
	
	\Needspace{4\baselineskip}
	\textbf{Schulz, Jan} and Mayerhoffer, Daniel (2021): ``Equal chances, unequal outcomes? Network-based evolutionary learning and the industrial dynamics of superstar firms'', \textit{Journal of Business Economics}, 91: 1357--1385.
	
	\Needspace{6\baselineskip}
	\section{Editorials}
	\Needspace{4\baselineskip}
	Mayerhoffer, Daniel, Olbrich, Hannah, \textbf{Schulz, Jan} and Tisch, Daria (2026): ``Perceptions of Wealth Inequality'', editorial for the special issue in \textit{Historical Social Research}. In print.
	
	\Needspace{4\baselineskip}
	\textbf{Schulz, Jan}, H\"otte, Kerstin and Mayerhoffer, Daniel (2024): ``Pluralist Economics in an Era of Polycrisis'', editorial for the special issue in the \emph{Review of Evolutionary Political Economy}, 5: 201--218.
	
	\section{Publications in Peer-Reviewed Edited Volumes} Mayerhoffer, Daniel and \textbf{Schulz, Jan} (2022): ``Marginalization and Misperception: Perceiving Gender and Racial Wage Gaps in Ego Networks''. In: Benito et al. (eds.): \textit{Complex Networks \& Their Applications X}. Springer.
	
	Mayerhoffer, Daniel and \textbf{Schulz, Jan} (2023): ``Redistribution, Social Segregation and Voting Information''. In: Sileno et al. (eds.) \textit{Proceedings of the 2nd Workshop on Agent-based Modeling and Policy-Making (AMPM 2022)}.  CEUR Workshop Proceedings.
	
	
	Mayerhoffer, Daniel, \textbf{Schulz, Jan} and Scheller, Simon (2021, in print): ``Multidimensionale Krisendiskurse in komplexen Gesellschaften''. In: Ott, M\"uhlnikel (eds.): \textit{Wege aus der Krise} (German-language edited volume).\\
	
	
	\section{Papers Under Review or Submitted} 
	Samartzidis, Lasare, Mundt, Philipp and \textbf{Schulz, Jan}: ``Production networks, input specificity, and the labor share''. BERG Working Paper No. 198. R\&R at \textit{Economic Modelling}.	
	
	
	Mayerhoffer, Daniel, Scheller, Simon, Schulz, Moritz and \textbf{Schulz, Jan}: ``Polarisation through Homophily. A Generative Algorithm''. Accepted abstract for the special issue ``Polarisation in Social Networks'' in \textit{Network Science} with R\&R at this journal.
	
	Van Langenhove, Christophe, Lorenz, Jan and \textbf{Schulz, Jan}: ``Wealth Taxation, Capital Gains Taxation and the Inequality--Mobility Trade-Off''. R\&R at the \textit{Journal of Economic Behavior and Organization}.
	
	\textbf{Schulz, Jan} and Ipsen, Leonhard: ``Poor Households and the Weight of Inflation''. FMM Working Paper No. 106. Under review at the \textit{Journal of Money, Credit and Banking}.
	Media coverage: ``Ein neuer Werkzeugkasten gegen die Inflation'', Makroskop (September 25, 2024).
	
	
	Olbrich, Hannah, Mayerhoffer, Daniel, Hartmann, Marie Lou, Gr\"otsch, Margitta and \textbf{Schulz, Jan}: ``Networks of Tunnels: Localised Perceptions and the Inequality-Wellbeing Link''. Invited contribution to a special issue in \textit{American Behavioral Scientist}; R\&R.
	
	\textbf{Schulz, Jan}, Memmen, Marvin and Ipsen, Leonhard: ``Sellers' inflation beyond collusion: Schumpeter meets Lerner''. Under review at the \textit{Journal of Economic Behavior and Organization}.
	
	\textbf{Schulz, Jan} and Weber, Jan D.: ``Power Laws in Socio-Economics''. Invited contribution to the Encyclopedia of Evolutionary and Institutional Economics.
	
	\textbf{Schulz, Jan}: ``Local and Global Inequality Rankings under Social Sampling''. Submitted.
	
	\textbf{Schulz, Jan}, Agoha, Caleb and Mayerhoffer, Daniel: ``Intersectional Privilege and Aggregate Wage Inequality''.  Submitted to \textit{The Journal of Economic Inequality}.
	\Needspace{5\baselineskip}
	\section{Work in Progress} 
	\Needspace{4\baselineskip}
	Rosina, Giulia, Mayerhoffer, Daniel, Noichl, Max and \textbf{Schulz, Jan}: ``The Eclipse of Inequality? A Bibliometric Analysis of Inequality Research in Economics''.
	
	\Needspace{4\baselineskip}
	Mayerhoffer, Daniel, Hartmann, Marie Lou, Olbrich, Hannah and \textbf{Schulz, Jan}: ``Narratives of Inequality and Zero-Sum Thinking: Evidence from a Representative Survey in Germany''. Data collection and preregistered main analyses completed; German adult sample of 1,575 and youth sample of 1,629.
	
	\Needspace{4\baselineskip}
	\textbf{Schulz, Jan}, Olbrich, Hannah E., Gr\"otsch, Margitta, Hartmann, Marie Lou and Mayerhoffer, Daniel M.: ``The Effect of Masculine versus Feminine Forms of Occupational Titles on Perceived Pay''. Working paper.
	
	\Needspace{4\baselineskip}
	\textbf{Schulz, Jan}, Fetai, Omer,	Gebhard, Anna and Milakovi\'c, Mishael: ``Aggregating Microeconomic Shocks in Production Networks''.
	
	\Needspace{4\baselineskip}
	\textbf{Schulz, Jan}: ``Spatial Social Structure in Germany, 1995--2009''. Fifteen geocoded nationwide telephone-directory cross-sections (roughly 30--35 million records per wave) for studying segregation, mobility, agglomeration, local shocks and political outcomes; preliminary analyses link local status composition to later voting, including AfD support.
	
	\Needspace{4\baselineskip}
	\textbf{Schulz, Jan} and Tsoukli, Xanthi: ``Class, Inequality, and the Social Organization of Military Risk in the First World War''.
	
	
	\Needspace{4\baselineskip}
	Dix, Jonas, Ipsen, Leonhard and \textbf{Schulz, Jan}: ``Decomposing the distributional effects of supply- and demand-side price shocks: An SVAR perspective''.
	
	\Needspace{4\baselineskip}
	Hartmann, Marie Lou, Olbrich, Hannah, Mayerhoffer, Daniel and \textbf{Schulz, Jan}: ``I don't see anything positive or neutral in it'': Adolescents' Lived Experiences, Perceptions and Understandings of Inequality.
	
	\Needspace{4\baselineskip}
	Olbrich, Hannah, Hartmann, Marie Lou, Mayerhoffer, Daniel  and \textbf{Schulz, Jan}:  ``Tied to Reality? - A Network-based Approach to Inequality Perceptions in the Classroom''.
	
	%	Ipsen, Leonhard, Rochowicz, Nils, Lindenmeyer, Adrian and \textbf{Schulz, Jan}: ``Friend- and Re-Shoring Decisions of Automotive Suppliers - A Network Perspective''.
	
	%Gr\"otsch, Margitta, Mayerhoffer, Daniel and \textbf{Schulz, Jan}: ``Fitness Landscapes and Structural Change''.
	
	
	%Mayerhoffer, Daniel and \textbf{Schulz, Jan}: ``Economic and Epistemic Privilege: An Agent-Based Exploration''.
	
	\section{Other Writings} 
	\textbf{Schulz, Jan} and Weber, Jan (2023): ``Is it simple? Is it complicated? No, it's complexity economics!'', \emph{Teaching Heterodox Economics Magazine}, 1(1): 32--33.
	
	Mayerhoffer, Daniel, \textbf{Schulz, Jan} and Tisch, Daria (2023): ``Fehlwahrnehmungen von \"Uberreichtum''. Text for the essay award on excessive wealth (``\"Uberreichtum'') by the Arbeiterkammer Wien.
	
	\section{Media Appearances} 
	B\"ockler Impuls 19/2021 ``Untersch\"atzter Reichtum'', Paper of the Week at Momentum Institut, interview with stepstone.de ``Ab wann ist man reich?'' (January 09, 2024), Makroskop (September 25, 2024), Podcast ``In der Wirtschaft'' (No. 61), Frankfurter Neue Presse (March 23, 2024), M\"unchener Merkur (March 31 and July 18, 2024), Handelsblatt (April 11, July 06 and September 09, 2025; February 25, 2026), RTL Info Geldgeil?! ``Geld regiert die Welt: Wie Superreiche eure Politik beeinflussen'' (April 28, 2026), Tagesspiegel (February 27, 2026)
	\newpage
	\section{Teaching} \textbf{Didactics Training}\\
	University Teacher Training, ZHD Bamberg/ProfiLehre Plus: 94 Arbeitseinheiten (70.5 hours); ``Zertifikat Hochschullehre der Bayerischen Universit\"aten''.
	
	\textsc{University of Bamberg} \\
	\textbf{Lecturer} \\
	since 2025: Lecture \textit{Makro\"okonomik I} (BA level; compulsory 6-ECTS module, approx. 200--300 students; major independent teaching responsibility)\\
	\Needspace{3\baselineskip}
	2018 -- 2025 (every term except Summer Terms 2023 and 2024): Tutorial for \textit{Makro\"okonomik I} (BA level)
	
	\textbf{Co-Instructor} \\
	Summer Term 2026: \textit{Constructing worlds! Die faszinierende Welt agentenbasierter Computersimulationen - Grundlagen und erste Anwendungen} (BA level, co-taught with Political Science)\\
	Summer Term 2024: The Political Economy of Inequality (co-taught with Political Science Institute, BA level)\\
	Summer Term 2019 -- : The Economics of Inequality (MA level; sole instructor with full course, supervision and assessment responsibility in Summer 2024)\\
	Summer Term 2020 -- : History of Economic Thought (BA level; sole instructor with full course, supervision and assessment responsibility in Summer 2024)\\
	Winter Term 2018 -- 2021: Current Issues in European Economic Policy (BA level)\\
	Winter Term 2020/21: Explanations in Economics and the Social Sciences (BA level, co-taught with Political Science Institute)\\
	Summer Term 2022: Computational Simulation for the Social Sciences (BA level, co-taught with Political Science Institute)\\
	since 2018: Co-Supervision of about 30 Bachelor and Master theses, primarily on economic inequality and international economics
	
	
	\Needspace{9\baselineskip}
	\textbf{Selected Official Teaching Evaluations} (EvaSys global indicator, 1--5 scale, 5 = best)\\
	Summer Term 2024: The Economics of Inequality, 5.0/5\\
	Summer Term 2026: The Economics of Inequality, 4.8/5\\
	Summer Term 2026: Constructing worlds! Computational Simulation, 4.8/5\\
	Winter Term 2023/24: \textit{Tutorium zu Makro\"okonomik I}, 4.8/5\\
	Summer Term 2024: History of Economic Thought, 4.7/5\\
	Summer Term 2022: Computational Simulation for the Social Sciences, 4.7/5\\
	Summer Term 2026: History of Economic Thought, 4.6/5
	
	\Needspace{5\baselineskip}
	\textsc{University of Bayreuth}\\
	\textbf{Teaching Assistant} \\
	Summer Term 2014: Social Choice Theory (BA and MA level)\\
	Winter Term 2012/13: Introduction to Philosophy (BA level)
	
	\textsc{Guest Lectures} \\
	Summer Term 2021: Lecture on Perceptions of Gender Inequality (for the course \textit{Striving for Gender Equality -- Contributions of Positive Political Theory}, University of Bamberg)\\
	
	Summer Term 2020: Lecture on the History of Neoliberalism (for the course \textit{US Economic History}, University of Utah)\\
	\newpage
	\section{Conference Presentations (including planned)}
	\textbf{2027}: Virtual Workshop in Historical Political Economy, `Class, Inequality, and the Social Organization of Military Risk in the First World War''.\\
	\textbf{2026:} EAEPE (Lausanne), ``Growing Differently: European Integration and Regional Cohesion''; 7th Pluralumn* Workshop, ``Class, Inequality, and the Social Organization of Military Risk in the First World War''; European PPE Network (Bayreuth), ``Excessive White Male Privilege Biases the Measurement of Intersectional Discrimination''.\\
	\textbf{2025:} Pluralumn* Workshop (Hamburg), ``Taxation, growth and the inequality-mobility nexus in multiplicative growth regimes''.\\
	\Needspace{9\baselineskip}
	\textbf{2024:} International Workshop on Economic Complexity and Macroeconomic Dynamics (Valencia), ``Aggregating Microeconomic Shocks in Production Networks''; Young Economists Conference (Vienna), ``Social Segregation, Misperceptions, and Emergent Cyclical Choice Patterns''; EAEPE (Bilbao), ``Poor Households and the Weight of Inflation''.\\
	\textbf{2023:} NetSci (Vienna), ``Regional value trees in Europe and Functional Distribution''.\\
	\textbf{2022:} AMPM; FMM; Young Economists Conference; EAEPE; BaGBeM Behavioral Macroeconomics Workshop; Networks in the Global World.\\
	\textbf{2021:} Complex Networks; FMM; Philosophy \& Economics Conference; DVPW Congress; EAEPE; WEHIA; ECINEQ; Networks Conference; PKES PhD Conference.\\
	\textbf{2020:} WICK; Young Economists Conference; EAEPE.\\
	\textbf{2019:} FMM; WEHIA; BaGBeM Behavioral Macroeconomics Workshop.\\
	\textbf{2018:} CFE; FMM; BaGBeM Behavioral Macroeconomics Workshop.\\
	
	\Needspace{5\baselineskip}
	\section{Invited Talks} 
	\Needspace{4\baselineskip}
	\textbf{Jul 2026} Lecture Series Plural Perspectives on Distribution: `` Wealth Distribution and Its Perception'', University of Cologne, Germany.
	
	\Needspace{4\baselineskip}
	\textbf{Jul 2024} Wealth in Data Science Summer School ``Some Aspects of Power Laws in Economics and the Social Sciences'' and group lead ``Redistribution in multiplicative growth regimes'', Bremen (Germany).
	
	\Needspace{4\baselineskip}
	\textbf{Apr 2022} NetPLACE (Networks, Phd Life And ComplExity Seminars) ``A Network-Based Explanation of Perceptions of Inequality and Privilege'', online.
	
	\Needspace{4\baselineskip}
	\textbf{Nov 2021} freie uni bamberg. ``Social Comparisons and Networked Perception'', Bamberg (Germany).
	
	\Needspace{4\baselineskip}
	\textbf{Aug 2021} Annual Conference of the International Institute of Public Finance 2021. Invited Talk to the panel ``The Super Rich'' (invitation by Dr. Isabel Mart\'inez), online.
	
	\Needspace{4\baselineskip}
	\textbf{May 2021} Friedrich-Ebert-Stiftung. Gesellschaftspolitisches Web-Seminar. Invited panel discussion: ``Schicksalsschlag, Katastrophe, Krise -- oder keiner Schlagzeile wert?'', online.
	
	
	\Needspace{4\baselineskip}
	\textbf{May 2021} Annual convention of Gr\"une Jugend Bamberg. ``Inequality, Consumption and the Critique of Consumption'', Bamberg (Germany).
	
	\Needspace{4\baselineskip}
	\textbf{Feb 2021} University of Utah Research Seminar. ``A Network-Based Explanation of Perceived Inequality'', online.
	
	\Needspace{4\baselineskip}
	\textbf{Nov 2020} Herbsttagung Netzwerk Plurale \"Okonomik. ``Agent-Based Modelling and Network Theory'', online.
	
	\newpage
	\section{Grants}
	\Needspace{5\baselineskip}
	\textbf{2026} FNK ``fres(c)h!'' grant, University of Bamberg, for ``Surnames, Socioeconomic Status, and the Burden of War: The Case of Imperial Germany in World War I'', with Xanthi Tsoukli (3,000 Euro total; 1,500 Euro notional share). \textit{Project:} Class composition of German First World War casualties using novel microdata with millions of observations; preliminary analyses show a negative mortality gradient by social status. A larger external grant application is planned.
	
	\Needspace{5\baselineskip}
	\textbf{2024} Bundesministerium f\"ur Bildung und Forschung (BMBF) grant ``e-lernen'' (project lead, 2024 -- 2026; 330,000 Euro total, 180,000 Euro Bamberg component under my leadership). \textit{Project:} Developing and evaluating an open e-learning course on economic inequality and studying inequality perceptions, social comparison and friendship networks among adolescents.
	
	\Needspace{5\baselineskip}
	\textbf{2023} Volkswagen Foundation (VolkswagenStiftung) grant for ``SPINE - Spring School on Perceptions of Wealth Inequality'' (project co-lead), University of Bamberg (March 2025, 78,300 Euro). \textit{Project:} Research incubator bringing together different groups working on inequality perceptions and leading to a special issue in \textit{Historical Social Research}.
	
	\Needspace{4\baselineskip}
	\textbf{2022} Event grant by the Institute for New Economic Thinking for the Third Pluralumn* Workshop (2,500 Dollar).
	
	\Needspace{5\baselineskip}
	\textbf{2021 -- 2022} Internal ``fres(c)h!'' grant by the University of Bamberg for the interdisciplinary project ``Social Comparisons and Inequality Perception in Homophilic Networks'' (3,000 Euro). \textit{Project:} Developing several simulation models on segregation and misperceptions of economic, gender, and racial inequality.
	
	\Needspace{4\baselineskip}
	\textbf{2018 -- 2020} PhD scholarship from the Hans-B\"ockler-Stiftung ($\approx$ 35,000 Euro).
	
	\Needspace{3\baselineskip}
	\textbf{2013} DAAD scholarship for semester abroad at the University of Graz.\\
	
	\Needspace{5\baselineskip}
	\section{Awards}
	\Needspace{4\baselineskip}
	\textbf{2026} ``Preis f\"ur gute Lehre'' of the Institut f\"ur Volkswirtschaftslehre, University of Bamberg (for Summer Term 2025).
	
	\Needspace{4\baselineskip}
	\textbf{2026} Nominee for the faculty-level ``Preis f\"ur gute Lehre'' 2025, Fakult\"at Sozial- und Wirtschaftswissenschaften, University of Bamberg.
	
	\Needspace{4\baselineskip}
	\textbf{2025} Poster award (3rd place) for the poster on ``Economic gender inequality, intersectionality and privilege: A participatory research agenda'', 100 Euro.
	
	\Needspace{4\baselineskip}
	\textbf{2024} Nominee for the faculty-level ``Preis f\"ur gute Lehre'' 2023, Fakult\"at Sozial- und Wirtschaftswissenschaften, University of Bamberg.
	
	\Needspace{4\baselineskip}
	\textbf{2023} Recognition award in the essay competition of the Arbeiterkammer Wien for ``Fehlwahrnehmungen von \"Uberreichtum'', with D. Mayerhoffer and D. Tisch, 1,500 Euro.
	
	\Needspace{4\baselineskip}
	\textbf{2022} Audience poster award for the poster on ``Perception and Privilege'' awarded by the Women's Affairs Office of the University of Bamberg, 250 Euro.
	
	\Needspace{4\baselineskip}
	\textbf{2022} Egon-Matzner-Preis f\"ur Sozio\"okonomie for the paper ``A Network-Based Explanation of Inequality Perceptions'' awarded by the Department for Public Finance and Infrastructure Policy (IFIP) of TU Vienna, 1,000 Euro.
	
	\Needspace{4\baselineskip}
	\textbf{2022} Herbert Simon Prize for the paper ``A Network Approach to Consumption'' awarded by the European Association for Evolutionary Political Economy (EAEPE).
	
	\Needspace{4\baselineskip}
	\textbf{2020} Nominee for the faculty-level ``Preis f\"ur gute Lehre'' 2019, Fakult\"at Sozial- und Wirtschaftswissenschaften, University of Bamberg.
	
	\section{Professional and Elected Roles}
	\textbf{since 2026} Associate Editor, \textit{SN Business \& Economics}.
	
	\textbf{since 2025} Member, Zentrum f\"ur Geschlechtersensible Forschung (ZGF) / Center for Gender-Sensitive Research, University of Bamberg; elected Steering Committee member since 2026.
	
	\section{Service} 
	
	\textbf{Reviewer and Editorial Activity}\\
	Associate Editor, \textit{SN Business \& Economics} (since 2026).\\
	
	Co-editor, Palgrave Macmillan (Springer Nature) conference proceedings volume \textit{Pluralism in Practice} (contract signed; forthcoming).\\
	
	Member of the Young Scholar Review Board of \textit{Social Science Research} (since 2025).
	
	Guest Editor for the Special Issue ``Perceptions of Wealth and Inequality'' of Historical Social Research (to be published in September 2026).
	
	Guest Editor for the Special Issue of the Review of Evolutionary Political Economy (REPE) ``Back to the future? Pluralist economics coming of age in an era of multiple crises'' (published in June 2024).
	
	\textbf{Refereeing}:	Journal of Business Economics (2020), Journal of Economic Interaction and Coordination (2021, 2022, 2025, 2026), Wege aus der Krise (edited volume, 2021), Springer Lecture Notes in Networks and Systems (2022), Acta Sociologica (2023), 28th Annual Meeting of the European Association of Young Economists (2023), Wealth in Data Science Summer School (2024), Journal of Post Keynesian Economics (2024), Social Science Research (2024, 2025, 2026), Historical Network Research Conference (2025), Review of Evolutionary Political Economy (2025), Scientific Reports (2025), International Journal of Empirical Economics (2025), Sage Open (2025), Quality \& Quantity (2025, 2026), Journal of Population Ageing (2026), Humanities and Social Sciences Communications (2026), British Journal of Social Psychology (2026), International Review of Economics and Finance (2026), Social Networks (2026), Journal of Economic Dynamics and Control (2026)
	
	
	\textbf{Conference Organization}\\
	Organizer for a special session at the Annual EAEPE Conference (2026) on ``Structural Inequality, Climate Policy, and Social Perception'', University of Lausanne.
	
	Scientific committee and organizing team for the 7th Pluralumn* workshop (2026), online.
	
	Scientific committee and organizing team for the 6th Pluralumn* workshop (2025), University of Hamburg and University of Flensburg.
	
	Scientific committee for the Historical Network Research Conference, Rio de Janeiro, 2025.
	
	Scientific committee and organizing team for the 5th Pluralumn* workshop: ``Empowering young minds'' (2024), University of Duisburg-Essen.
	
	Scientific committee and organizing team of the satellite workshop ``From Micro to Macro Via Network Interaction'' of the NetSci2023 conference, Vienna (Austria).
	
	
	Scientific committee and organizing team for the 4th Pluralumn* workshop: ``Pluralist economics beyond anything goes'' (2023), University of Bamberg. 
	
	
	Scientific committee and organizing team for the 3rd Pluralumn* workshop: ``Back to the future? Pluralist economics coming of age in an era of multiple crises'' (2022), University of Bielefeld. 
	
	Organizing team for the Second BaGBeM Behavioral Macroeconomics Workshop (2019), University of Bamberg.
	
	\textbf{Pluralist Economics}\\
	Founding Member of the Bamberg Group for Pluralist Economics
	(organization and speaker at about 15 lectures and seminars for the Bamberg Group for Pluralist Economics  (2015 -- 2023), Editor at the online learning platform ``Exploring Economics'' (2020 -- 2022).\\
	
	
	\section{Selected Summer Schools and Workshops}
	\textbf{2024:} Modelling Shock Propagation and Systemic Risk in Supply Networks, Bamberg Doctoral Research Training Group.\\
	\textbf{2019:} Winter School, Complexity Science Hub Vienna, Obergurgl (Austria).\\
	\textbf{2018:} Oxford Summer School on Economic Networks, Oxford (UK).\\
	
	\section{Memberships}
	Arbeitskreis Kritische Unternehmens- und Industriegeschichte (AKKU e.V.), Arbeitskreis Feministische Politische \"Okonomie, German Network for New Economic Dynamics, European Association for Evolutionary Political Economy (EAEPE), Netzwerk Plurale \"Okonomik, Bamberg Group for Pluralist Economics, Post-Keynesian Economics Society, Verein f\"ur Socialpolitik (German Economic Association).\\
	
	\Needspace{6\baselineskip}
	\section{Civic and University Service}
	\textbf{2015 -- 2017} Elected member of the Student Parliament, University of Bamberg, and its Gender Equality Working Group; spokesperson for social affairs in the Student Council.\\
	\textbf{2009 -- 2011} Chair of the Youth Parliament, Lauffen a.N.\\
	
	\enlargethispage{3\baselineskip}
	\hfill CV last updated 09/2026.
\end{resume}
\end{document}